% 5.4 空间群表示表的使用示例
\section{空间群表示表的使用示例}

立方密堆积与金刚石结构（$Fm3m\ (O^{5}_{h})$ 与 $Fd3m\ (O^{7}_{h})$）的空间群表示已在 3.8 节作为推导空间群表示的理论之例导出。现在我们给出一个例子，说明如何使用表 5.7 与 5.8 为 230 个空间群中的任何一个检索类似的信息。我们留给读者一个附加练习：验证 3.8 节讨论过的 $Fm3m\ (O^{5}_{h})$ 与 $Fd3m\ (O^{7}_{h})$ 的空间群表示如何能从表 5.7 与 5.8 得到。

\begin{figure}[H]
\centering
\includegraphics[width=0.62\textwidth]{fig5_2_mu.jpg}
\caption*{\textbf{图 5.2}\ $\Gamma^{f}_{c}$ 的布里渊区。$\Gamma=(000)$；$X=(\frac{1}{2}0\frac{1}{2})$；$L=(\frac{1}{2}\frac{1}{2}\frac{1}{2})$；$W=(\frac{1}{2}\frac{1}{4}\frac{3}{4})$。$X$ 与 $U$ 虽常被称为对称点，但按定义 3.3.1 它们并不算对称点，因为它们并不比相邻的线 $\Delta$ 与 $S$ 具有更多对称性。}
\end{figure}

空间群 $F\bar{4}3c\ (T^{5}_{d})$ 在表 5.7 的注中已多次用作示例，这里用它进一步演示该表的用法。$F\bar{4}3c\ (T^{5}_{d})$ 是关联于点群 $\bar{4}3m\ (T_{d})$、建立在面心立方布拉维格子 $\Gamma^{f}_{c}$ 上的立方空间群，由表 3.1 其基本矢量为
\begin{equation*}
\left\{
\begin{array}{l}
\mathbf{t}_{1} = \frac{1}{2}(0, a, a),\\
\mathbf{t}_{2} = \frac{1}{2}(a, 0, a),\\
\mathbf{t}_{3} = \frac{1}{2}(a, a, 0),
\end{array}
\right.
\tag{5.4.1}
\end{equation*}
而
\begin{equation*}
\left\{
\begin{array}{l}
\mathbf{g}_{1} = (2\pi/a)(-1, 1, 1),\\
\mathbf{g}_{2} = (2\pi/a)(1, -1, 1),\\
\mathbf{g}_{3} = (2\pi/a)(1, 1, -1).
\end{array}
\right.
\tag{5.4.2}
\end{equation*}
因此倒格子是体心立方格子（见表 3.3），其布里渊区示于图 3.14 并复制为图 5.2。由表 3.6（$\Gamma^{f}_{c}$ 名下）或表 5.7，该布里渊区中特殊的对称点是 $\Gamma$、$X$、$L$ 与 $W$；一些作者把 $K$ 也当作对称点，但如图 3.14 的题注所指出的，它不符合我们关于对称点的定义（见定义 3.3.1）。类似地，对称线为 $\Delta$、$\Lambda$、$\Sigma$、$S$、$Z$ 与 $Q$。对称面不予考虑。各对称点和对称线的 $\mathbf{k}$ 矢量在表 3.6 中以 $\mathbf{g}_{1}, \mathbf{g}_{2}, \mathbf{g}_{3}$ 的单位给出：
\begin{center}
\begin{tabular}{lc}
点或线 & $\mathbf{k}$\\
\midrule
$\Gamma$ & $(000)$\\
$X$ & $(\frac{1}{2}0\frac{1}{2})$\\
$L$ & $(\frac{1}{2}\frac{1}{2}\frac{1}{2})$\\
$W$ & $(\frac{1}{2}\frac{1}{4}\frac{3}{4})$\\
$\Delta$ & $(\alpha 0\alpha)$\\
$\Lambda$ & $(\alpha\alpha\alpha)$\\
$\Sigma$ & $(\alpha, \alpha, 2\alpha)$\\
$S$ & $(\frac{1}{2}+\alpha, 2\alpha, \frac{1}{2}+\alpha)$\\
$Z$ & $(\frac{1}{2}, \alpha, \frac{1}{2}+\alpha)$\\
$Q$ & $(\frac{1}{2}, \frac{1}{2}-\alpha, \frac{1}{2}+\alpha)$
\end{tabular}
\end{center}

我们逐一考虑各个对称点。

\noindent\textbf{$\Gamma$：$\mathbf{k} = (000)$。}

${}^{H}\mathbf{G}^{\Gamma}$ 是 $\mathbf{G}^{7}_{24}$，它同构于点群 $\bar{4}3m\ (T_{d})$，其生成元为
\begin{equation*}
P = \{C^{-}_{31} \mid 000\}, \qquad Q = \{C_{2z} \mid 000\},
\end{equation*}
\begin{equation*}
R = \{C_{2x} \mid 000\}, \qquad S = \{\sigma_{da} \mid {\textstyle\frac{1}{2}}{\textstyle\frac{1}{2}}{\textstyle\frac{1}{2}}\}.
\end{equation*}
五个类（见表 5.1）为
\begin{center}
\begin{tabular}{l}
$C_{1} = \{E \mid 000\}$；\\
$C_{2} = \{C_{2x} \mid 000\}, \{C_{2y} \mid 000\}, \{C_{2z} \mid 000\}$；\\
$C_{3} = \{S^{-}_{4x} \mid {\frac{1}{4}}{\frac{3}{4}}{\frac{1}{4}}\}$ 及其余五个 $S^{\pm}_{4m}$ 项；\\
$C_{4} = \{\sigma_{da} \mid {\frac{1}{4}}{\frac{3}{4}}{\frac{1}{4}}\}$ 及其余五个 $\sigma_{dm}$ 项；\\
$C_{5} = \{C^{-}_{31} \mid 000\}, \dots, \{C^{+}_{34} \mid 000\}$（八个 $C^{\pm}_{3j}$）。
\end{tabular}
\end{center}
要辨认 $\mathbf{G}^{7}_{24}$ 的哪个元素对应于哪个空间群元素相当费事，必须直接相乘。见表 5.7 的注 (iv)，那里证明了 $\mathbf{G}^{7}_{24}$ 的 $P^{2}RS$ 对应于 $\{S^{-}_{4y} \mid {\frac{1}{2}}\,{-\frac{3}{4}}\,{\frac{1}{4}}\}$，或平移至 $\mathbf{F}$ 后为 $\{S^{-}_{4y} \mid {\frac{1}{2}}{\frac{1}{4}}{\frac{1}{4}}\}$。如果要求的是简并表示的实际矩阵代表而不仅是特征标，这种认定当然是必须做的。于是 ${}^{H}\mathbf{G}^{\Gamma}$ 的特征标表可由表 5.1 的 $\mathbf{G}^{7}_{24}$ 条目查得，只是表示不标记为 $R_{1}, \dots, R_{5}$，而按表 5.8 中 $\mathbf{G}^{7}_{24}$ 下的条目“a”标记：
\begin{center}
\begin{tabular}{lccccc}
 & $R_{1}$ & $R_{2}$ & $R_{3}$ & $R_{4}$ & $R_{5}$\\
a & $A_{1}$ $\Gamma_{1}$ & $A_{2}$ $\Gamma_{2}$ & $E$ $\Gamma_{3}$ & $T_{1}$ $\Gamma_{4}$ & $T_{2}$ $\Gamma_{5}$
\end{tabular}
\end{center}
由表 5.7 可见这些空间群表示全为第一种，于是 ${}^{H}\mathbf{G}^{\Gamma}$ 的特征标表为（右列给出相应空间群表示的实性）：
\begin{center}
\begin{tabular}{lcccccc}
 & $C_{1}$ & $C_{2}$ & $C_{3}$ & $C_{4}$ & $C_{5}$ & \\
$\Gamma_{1}\ A_{1}$ & 1 & 1 & 1 & 1 & 1 & 1\\
$\Gamma_{2}\ A_{2}$ & 1 & 1 & $-1$ & 1 & 1 & 1\\
$\Gamma_{3}\ E$ & 2 & 2 & 0 & 0 & $-1$ & 1\\
$\Gamma_{4}\ T_{1}$ & 3 & $-1$ & 1 & $-1$ & 0 & 1\\
$\Gamma_{5}\ T_{2}$ & 3 & $-1$ & $-1$ & 1 & 0 & 1
\end{tabular}
\end{center}
若计及时间反演对称，也没有额外简并。

\noindent\textbf{$X$：$\mathbf{k} = (\frac{1}{2}0\frac{1}{2})$。}

${}^{H}\mathbf{G}^{X}$ 是 $\mathbf{G}^{4}_{8}$ 与 $\mathbf{T}_{2}$ 的直积，$\mathbf{G}^{4}_{8}$ 的生成元为
\begin{equation*}
P = \left\{S^{+}_{4y} \mid {\textstyle\frac{1}{2}}{\textstyle\frac{1}{2}}{\textstyle\frac{1}{2}}\right\}, \qquad Q = \{C_{2x} \mid 000\},
\end{equation*}
$\mathbf{T}_{2}$ 的元素为 $\{E \mid \mathbf{0}\}$ 与 $\{E \mid \mathbf{t}_{s}\}$（$s = 1$ 或 3）。于是 $\mathbf{G}^{4}_{8}$ 的元素为
\begin{center}
\begin{tabular}{llll}
$E$ & $\{E \mid 000\}$ & $Q$ & $\{C_{2x} \mid 000\}$\\
$P$ & $\{S^{+}_{4y} \mid {\frac{1}{2}}{\frac{1}{2}}{\frac{1}{2}}\}$ & $PQ$ & $\{C_{2c} \mid {\frac{1}{2}}{\frac{1}{2}}{\frac{1}{2}}\}$\\
$P^{2}$ & $\{C_{2y} \mid 000\}$ & $P^{2}Q$ & $\{C_{2z} \mid 000\}$\\
$P^{3}$ & $\{S^{-}_{4y} \mid {\frac{1}{2}}{\frac{1}{2}}{\frac{1}{2}}\}$ & $P^{3}Q$ & $\{C_{2e} \mid {\frac{1}{2}}{\frac{1}{2}}{\frac{1}{2}}\}$
\end{tabular}
\end{center}
由表 5.1 它们可归入类 $C_{1}, C_{2}, \dots, C_{5}$。五个表示按表 5.8 标记为 $A_{1}, A_{2}, B_{1}, B_{2}$ 与 $E$（即 $X_{1}, X_{2}, X_{3}, X_{4}$ 与 $X_{5}$）。于是 ${}^{H}\mathbf{G}^{X}$ 的特征标表为（右列给出实性）：
\begin{center}
\begin{tabular}{lccccccc}
 & & $\{E \mid 000\}$ & $\{C_{2y} \mid 000\}$ & $\begin{array}{c}\{S^{+}_{4y} \mid {\frac{1}{4}}{\frac{3}{4}}{\frac{1}{4}}\}\\ \{S^{-}_{4y} \mid {\frac{1}{4}}{\frac{3}{4}}{\frac{1}{4}}\}\end{array}$ & $\begin{array}{c}\{C_{2x} \mid 000\}\\ \{C_{2z} \mid 000\}\end{array}$ & $\begin{array}{c}\{C_{2c} \mid {\frac{1}{4}}{\frac{3}{4}}{\frac{1}{4}}\}\\ \{C_{2e} \mid {\frac{1}{4}}{\frac{3}{4}}{\frac{1}{4}}\}\end{array}$ & \\
$X_{1}$ & $A_{1}$ & 1 & 1 & 1 & 1 & 1 & 1\\
$X_{2}$ & $A_{2}$ & 1 & 1 & 1 & $-1$ & $-1$ & 1\\
$X_{3}$ & $B_{1}$ & 1 & 1 & $-1$ & 1 & $-1$ & 1\\
$X_{4}$ & $B_{2}$ & 1 & 1 & $-1$ & $-1$ & 1 & 1\\
$X_{5}$ & $E$ & 2 & $-2$ & 0 & 0 & 0 & 1
\end{tabular}
\end{center}
若计及时间反演对称，将没有额外简并。若需要表示 $E\ (X_{5})$ 的矩阵，它们也可由表 5.1 中 $\mathbf{G}^{4}_{8}$ 的 $R_{5}$ 得到。

\noindent\textbf{$L$：$\mathbf{k} = (\frac{1}{2}\frac{1}{2}\frac{1}{2})$。}

${}^{H}\mathbf{G}^{L}$ 是 $\mathbf{G}^{4}_{12}$，其生成元为
\begin{equation*}
P = \{C^{-}_{31} \mid 001\}, \qquad Q = \{\sigma_{db} \mid {\textstyle\frac{1}{2}}{\textstyle\frac{1}{2}}{\textstyle\frac{1}{2}}\},
\end{equation*}
而 ${}^{H}\mathbf{G}^{L}$ 的十二个成员为
\begin{center}
\begin{tabular}{llll}
$E$ & $\{E \mid 000\}$ & $Q$ & $\{\sigma_{db} \mid {\frac{1}{2}}{\frac{1}{2}}{\frac{1}{2}}\}$\\
$P$ & $\{C^{-}_{31} \mid 001\}$ & $PQ$ & $\{\sigma_{df} \mid {\frac{1}{2}}{\frac{1}{2}}{-\frac{1}{2}}\}$\\
$P^{2}$ & $\{C^{+}_{31} \mid 000\}$ & $P^{2}Q$ & $\{\sigma_{de} \mid {\frac{1}{2}}{\frac{1}{2}}{\frac{1}{2}}\}$\\
$P^{3}$ & $\{E \mid 001\}$ & $P^{3}Q$ & $\{\sigma_{db} \mid {\frac{1}{2}}{\frac{1}{2}}{-\frac{1}{2}}\}$\\
$P^{4}$ & $\{C^{-}_{31} \mid 000\}$ & $P^{4}Q$ & $\{\sigma_{df} \mid {\frac{1}{2}}{\frac{1}{2}}{\frac{1}{2}}\}$\\
$P^{5}$ & $\{C^{+}_{31} \mid 001\}$ & $P^{5}Q$ & $\{\sigma_{de} \mid {\frac{1}{2}}{\frac{1}{2}}{-\frac{1}{2}}\}$
\end{tabular}
\end{center}
它们可归入 $\mathbf{G}^{4}_{12}$ 的类，从而构造出 ${}^{H}\mathbf{G}^{L}$ 的特征标表：
\begin{center}
\begin{tabular}{lccccccc}
 & $\{E \mid 000\}$ & $\{E \mid 001\}$ & $\begin{array}{c}\{C^{-}_{31} \mid 001\}\\ \{C^{+}_{31} \mid 001\}\end{array}$ & $\begin{array}{c}\{C^{+}_{31} \mid 000\}\\ \{C^{-}_{31} \mid 000\}\end{array}$ & $\begin{array}{c}\{\sigma_{db} \mid {\frac{1}{2}}{\frac{1}{2}}{\frac{1}{2}}\}\\ \{\sigma_{de}\},\{\sigma_{df}\}\end{array}$ & $\begin{array}{c}\{\sigma_{db} \mid {\frac{1}{2}}{\frac{1}{2}}{-\frac{1}{2}}\}\\ \{\sigma_{de}\},\{\sigma_{df}\}\end{array}$ & \\
$L_{1}$\ \ ${}^{2}E$ & 1 & $-1$ & $-1$ & 1 & $i$ & $-i$ & 3\\
$L_{2}$\ \ ${}^{1}E$ & 1 & $-1$ & $-1$ & 1 & $-i$ & $i$ & 3\\
$L_{3}$\ \ $E$ & 2 & $-2$ & 1 & $-1$ & 0 & 0 & 2
\end{tabular}
\end{center}
$E\ (L_{3})$ 的矩阵可由表 5.1 中 $\mathbf{G}^{4}_{12}$ 的 $R_{6}$ 得到。最右列的实性表明：若计及时间反演对称，${}^{1}E$ 与 ${}^{2}E$（$L_{2}$ 与 $L_{1}$）之间出现额外简并，属于表示 $E\ (L_{3})$ 的两组不同基函数之间也出现额外简并。

\noindent\textbf{$W$：$\mathbf{k} = (\frac{1}{2}\frac{1}{4}\frac{3}{4})$。}

${}^{H}\mathbf{G}^{W}$ 是 $\mathbf{G}^{1}_{4}$ 与 $\mathbf{T}_{4}$ 的直积，其中 $\mathbf{T}_{4}$ 由 $\{E \mid \mathbf{0}\}$、$\{E \mid \mathbf{t}_{2}\}$、$\{E \mid 2\mathbf{t}_{2}\}$、$\{E \mid 3\mathbf{t}_{2}\}$ 组成，而 $\mathbf{G}^{1}_{4}$ 的生成元为 $P = \{S^{+}_{4x} \mid {\frac{1}{4}}{\frac{1}{2}}{\frac{1}{4}}\}$。$\mathbf{G}^{1}_{4}$ 的每个元素自成一类，于是 ${}^{H}\mathbf{G}^{W}$ 的特征标表为
\begin{center}
\begin{tabular}{lccccc}
 & $\{E \mid 000\}$ & $\{S^{+}_{4x} \mid {\frac{1}{4}}{\frac{1}{2}}{\frac{1}{4}}\}$ & $\{C_{2x} \mid 000\}$ & $\{S^{-}_{4x} \mid {\frac{1}{4}}{\frac{1}{2}}{\frac{1}{4}}\}$ & \\
$W_{1}\ A$ & 1 & 1 & 1 & 1 & 3\\
$W_{2}\ {}^{2}E$ & 1 & $i$ & $-1$ & $-i$ & 3\\
$W_{3}\ B$ & 1 & $-1$ & 1 & $-1$ & 3\\
$W_{4}\ {}^{1}E$ & 1 & $-i$ & $-1$ & $i$ & 3
\end{tabular}
\end{center}
最右列的实性表明：若计及时间反演对称，属于不同表示的本征函数对之间出现额外简并；具体地说，$A\ (W_{1})$ 与 $B\ (W_{3})$ 变为简并，${}^{1}E\ (W_{4})$ 与 ${}^{2}E\ (W_{2})$ 变为简并（见 7.6 与 7.7 节）。

至此本空间群的对称点考虑完毕，接下来考虑对称线，这要用到 3.7 节描述的射影表示理论。对每条对称线，我们先辨认中央延拓群 $\overline{\mathbf{G}}^{\mathbf{k}*}$，再由它得到小群 $\mathbf{G}^{\mathbf{k}}$ 的表示。$\overline{\mathbf{G}}^{\mathbf{k}*}$ 的元素形如 $(H_{j}, \alpha)$，其中 $H_{j}$ 是某个空间群元素的点群操作部分，$\alpha$ 是式 (3.7.27) 定义的循环群 $\mathbf{Z}_{g}$ 的一个小整数。群 $\overline{\mathbf{G}}^{\mathbf{k}*}$ 在表 5.7 中指明，其生成元 $P, Q, \dots$ 依次以 $(H_{j}, \alpha)$ 的形式给出。于是 $\overline{\mathbf{G}}^{\mathbf{k}*}$ 的特征标表与矩阵代表可由表 5.1 辨认，而 $\mathbf{G}^{\mathbf{k}}$ 的表示对每个元素 $\{R \mid \mathbf{v}\}$ 由 $(R, 0)$ 的矩阵代表乘以因子 $\exp(-\mathrm{i}\mathbf{k} \cdot \mathbf{v})$ 得到（见式 (3.8.8) 与 (3.8.9)）。表示的标记可查表 5.7 中相应线条目所指出的表 5.8 的部分。注意，当 $\overline{\mathbf{G}}^{\mathbf{k}*}$ 的表示与 $(E, 1)$ 生成的循环群 $\mathbf{Z}_{g}$ 的允许表示不相容时，它们被略去（见式 (3.7.33)）。还应注意，一般 $\mathbf{G}^{\mathbf{k}}$ 的类结构可能与 $\overline{\mathbf{G}}^{\mathbf{k}*}$ 的不同，而且由于 $\mathbf{G}^{\mathbf{k}}$ 无限，确定 $\mathbf{G}^{\mathbf{k}}$ 的类结构可能相当复杂。相应空间群表示的实性（从而计及时间反演对称后的额外简并）也可由表 5.7 得到。下面以所选空间群 $F\bar{4}3c\ (T^{5}_{d})$ 的对称线为例说明。

以 $\Delta$ 线为例：中央延拓群 $\overline{\mathbf{G}}^{\Delta*}$ 是群 $\mathbf{G}^{2}_{4}$，其生成元为
\begin{equation*}
P = (C_{2y}, 0), \qquad Q = (\sigma_{dc}, 0).
\end{equation*}
由于 $\mathbf{G}^{2}_{4}$ 的阶为 4（与含 $E, C_{2y}, \sigma_{dc}, \sigma_{de}$ 的点群同阶），显然 $g$ 必为 1，于是 $\overline{\mathbf{G}}^{\Delta*}$ 的四个元素按类为：$C_{1}$：$(E, 0)$；$C_{2}$：$(C_{2y}, 0)$；$C_{3}$：$(\sigma_{dc}, 0)$；$C_{4}$：$(\sigma_{de}, 0)$。由表 5.1 得 $\overline{\mathbf{G}}^{\Delta*}$ 的特征标表，并用表 5.8 给出表示的标记。$\mathbf{G}^{\Delta}$ 的小表示即上表每项乘以因子 $\exp(-\mathrm{i}\mathbf{k} \cdot \mathbf{v})$：
\begin{center}
\begin{tabular}{lccccc}
 & $\{E \mid 000\}$ & $\{C_{2y} \mid 000\}$ & $\{\sigma_{dc} \mid {\frac{1}{2}}{\frac{1}{2}}{\frac{1}{2}}\}$ & $\{\sigma_{de} \mid {\frac{1}{2}}{\frac{1}{2}}{\frac{1}{2}}\}$ & \\
$\Delta_{1}\ A_{1}$ & 1 & 1 & $\Lambda$ & $\Lambda$ & 1\\
$\Delta_{3}\ A_{2}$ & 1 & 1 & $-\Lambda$ & $-\Lambda$ & 1\\
$\Delta_{2}\ B_{1}$ & 1 & $-1$ & $\Lambda$ & $-\Lambda$ & 3\\
$\Delta_{4}\ B_{2}$ & 1 & $-1$ & $-\Lambda$ & $\Lambda$ & 3
\end{tabular}
\end{center}
其中 $\Lambda = \exp(-2\pi\mathrm{i}\alpha)$。计及时间反演对称后，$A_{1}\ (\Delta_{1})$ 与 $A_{2}\ (\Delta_{3})$ 没有额外简并，但 $B_{1}\ (\Delta_{2})$ 与 $B_{2}\ (\Delta_{4})$ 变为简并。

\noindent（中译者注：原书随后对 $\Lambda, \Sigma, S, Z, Q$ 各线按同一 recipe 逐一处理（原书第 402--405 页），即辨认中央延拓群 $\overline{\mathbf{G}}^{\mathbf{k}*}$（表 5.1 + 表 5.7）并乘以相因子 $\exp(-\mathrm{i}\mathbf{k} \cdot \mathbf{v})$，此处从略；结果可直接由表 5.7 读出。）

空间群表示的用途现已扩展到晶体 lattice 正则振动的标记，即声子色散曲线的标记（Johnson and Loudon 1964；综述见 Maradudin and Vosko (1968) 与 Warren (1968)），也扩展到磁性晶体中自旋波色散关系的标记（Dimmock and Wheeler 1962\textit{b}, Joshua and Cracknell 1969, Loudon 1968）。最初声子色散曲线的实验测量只限于元素立方材料的高对称方向；在这种情形正则模可以相当充分地分类并标记为纯纵声学（LA）、横声学（TA）、纵光学（LO）或横光学（TO）。然而这个简单方案不适用于低对称方向或更复杂的晶体，所以现在正则模通常用空间群表示来标记。正如分子的正则振动可以群论地研究并归属于该分子点群的各表示（例如见 Cotton (1963), Cracknell (1968\textit{b}), Leech and Newman (1969)），具有给定波矢 $\mathbf{k}$ 的晶体正则振动也可以归属于该晶体空间群 $\mathbf{G}$ 的各表示。做法是研究晶胞中原子的一般位移 $\mathbf{x}_{i}$ 在空间群 $\mathbf{G}$ 操作下的变换性质。这样得到以 $\mathbf{x}_{i}$ 为基的表示 $\Gamma$，它一般可约。$\mathbf{k}$ 处正则模所属的 $\mathbf{G}^{\mathbf{k}}$ 的小表示于是可由 $\Gamma$ 的约化得到，小表示的标记随后用来标记布里渊区中 $\mathbf{k}$ 矢量的声子色散关系。另一种做法是：通过研究产生与湮灭算符 $a^{\dagger}_{\mathbf{k}}$ 与 $a_{\mathbf{k}}$ 在 $\mathbf{G}$ 操作下的变换性质、并把它们归属到 $\mathbf{G}^{\mathbf{k}}$ 的相应小表示，来确定波矢为 $\mathbf{k}$ 的正则模的对称性质；$a^{\dagger}_{\mathbf{k}}$ 与 $a_{\mathbf{k}}$ 是产生与湮灭波矢为 $\mathbf{k}$ 的声子的算符，其形式例如见 Kittel (1963)。磁性晶体自旋波色散关系按空间群表示的标记，可以通过研究类似于 $a^{\dagger}_{\mathbf{k}}$、$a_{\mathbf{k}}$ 的、产生与湮灭磁振子的相应算符在 $\mathbf{G}$ 操作下的变换性质来确定（Loudon 1968）。对某些特定空间群对称适配函数的确定，可见 Cornwell (1969) 一书的附录 3。在相对论性能带结构计算中必须使用双值空间群表示，这将在第六章讨论。
